\section{引言}

人工智能正在改变金融风险的识别、计量与防控方式，也在改变资产管理的实现路径。
大模型、图神经网络与强化学习等方法与复杂网络、多主体仿真等系统工程方法相互融合，
为市场风险与流动性风险的计量提供了新的工具\cite{hong2021,suzhi2017,wangfang2020}。
在这一背景下，多智能体量化投资的核心任务是把多个预测信号合成为可执行的组合。
机构投资者与量化管理人通常同时维护动量、反转、低波动、换手、估值与情绪等多个信号，
合成方式有两种。
静态合成以固定权重平均各信号；状态自适应合成（本文称门控组合）依据市场状态、
各信号近期表现与风险暴露动态调整信号权重。
随着多智能体系统与深度强化学习组合管理框架的流行\cite{jiang2017,liu2020,li2024}，
门控组合已成为多智能体量化投资中的默认架构。这一架构隐含地假设状态自适应优于静态合成，
但该假设缺乏明确的经济学条件。已有研究更多关注架构设计本身——增加智能体数量、
引入记忆机制、构造更复杂的门控网络——较少回答三个更基本的问题：

其一，条件问题：在什么条件下，状态自适应门控能带来扣除交易成本后的可测超额收益？
其二，来源问题：门控收益来自对不同状态下信号收益水平的预测（收益择时），
还是来自对不同状态下信号风险结构的配置（风险配置）？
其三，识别问题：给定样本的时间长度与横截面宽度，这一收益能否被统计地识别？

三个问题的答案决定风险管理资源的配置方向。若门控收益来自收益择时，
其识别依赖时间维度，在市场结构变化时容易失效；若收益来自风险配置，
低波动与风险平价等更简单的风险结构策略可能是更高效的实现方式。
对资产管理机构而言，判断门控是否值得做决定了模型容量、数据采购与交易系统的投入；
对监管者而言，人工智能投顾与量化资管中的「自适应」「智能」等业绩宣称需要可复核的判定标准。

本文研究多智能体量化投资策略的价值边界。核心判断是：复杂门控并不必然优于简单平均。
门控只有在信号的预测能力随市场状态出现足够大的差异（互补性），
且这一差异超过状态信息的获取成本与交易成本时才有价值；
否则，复杂的门控只是在为无法利用的状态差异支付估计与交易成本。
这一判断把「何时使用门控」从架构偏好转变为成本收益比较，
并把门控收益的性质界定为信息租金：它与信息经济学中「信息有价、
但价格不能长期高于获取成本」的逻辑同源\cite{grossman1980}。

在长仓满仓这一标准约束下（$w_t\in\{w\in\R^N_+:\bone^\top w=1\}$），组合收益可分解为
\begin{equation}
r_{p,t+1}=w_t^\top r_{t+1}=r_{ew,t+1}+\wtilde w_t^\top r_{t+1},
\qquad \wtilde w_t=w_t-w^{ew},\quad w^{ew}=\bone/N,
\label{eq:intro-decomp}
\end{equation}
其中 $r_{ew,t+1}$ 为等权组合收益。由 $\bone^\top\wtilde w_t=0$，组合的相对表现完全由
主动权重 $\wtilde w_t$ 决定，而 $r_{ew,t+1}$ 与决策无关。该分解的直接含义是：
以绝对收益为训练目标时，目标中与决策无关的基准项不提供梯度信息，
唯一较强的梯度来自换手惩罚，优化因此收敛到「不交易」这一平凡解。
命题 \ref{prop:degeneracy} 给出退化发生的精确参数条件。

该诊断有直接的实证对应物。在一个以十只美国 ETF 为样本、以绝对收益为目标的多智能体门控模型中，
平均换手率仅为 0.0029，与静态持有水平相当；组合夏普比率（0.940）与等权基准（0.941）
几乎一致，配对检验的 95\% 区间跨越零、Holm 校正后 $p=1.000$。
这一「自适应门控」实际上退化为静态等权，其零结果不是实验失败，而是目标函数设计的产物。
该现象影响文献中一类结论的可信度：在满仓长仓约束下报告「自适应优势」的研究，
若不同时报告组合相对基准的主动暴露，读者无法区分机制奏效与机制退化。

把训练目标改为主动收益即可消除退化：式 \eqref{eq:intro-decomp} 中被消去的基准项
不再主导梯度，门控具备识别信号—状态互补性的条件。

本文建立信号—状态互补性框架。门控相对最优静态合成的价值等于互补性收益
减去估计误差成本与换手成本，门控占优的充要条件是一个具有闭式解的阈值不等式
\begin{equation}
\hat\Gamma>\Theta+2\lambda\kappa\,\mathcal{C}(\hat g),
\label{eq:intro-threshold}
\end{equation}
其中 $\hat\Gamma$ 为互补性统计量，$\Theta$ 为估计噪声底线，$\kappa$ 为单边交易成本。
互补性可分解为一阶（条件均值）与二阶（条件波动与相关结构）两条渠道，
两者的噪声底线分别以 $O(1/T)$ 与 $O(1/(NT))$ 收敛；制度的马尔可夫持续性通过几何乘数
放大门控的信息价值；在容量约束下净价值关于制度数呈单峰，最优制度数内点存在；
横截面宽度不足时，互补性在统计上不可识别。

实证部分使用 2006 年 1 月至 2026 年 7 月中国 A 股月度点时面板（868 只证券、
247 个月度形成期、月均进入股票池 438 只），以 2014 年至 2025 年 12 个年度滚动窗口
进行样本外检验。

门控组合的表现与来源。门控组合的年化主动收益为 3.77\%，主动信息比率为 0.80，
相对等权的夏普差值为 $+0.115$（95\% 区间 $[-0.006,0.225]$，Holm 校正后 $p=0.386$），
经济量级明显而统计上不显著。按定理 \ref{thm:threshold}，
这一结果与两条渠道的互补性检验互相印证：一阶与二阶互补性统计量分别为
$3.93\times10^{-4}$ 与 $6.34\times10^{-4}$，对应噪声底线 $4.17\times10^{-4}$ 与
$6.87\times10^{-4}$，$p$ 值为 0.060 与 0.065。安慰剂检验把门控主动收益由 3.77\% 压至
0.004\%，说明收益依赖制度信息；而允许使用制度内真实条件期望收益的 Oracle 门控
（夏普 0.562）未跑赢静态合成（0.546），说明可利用的一阶互补性有限。
门控收益主要来自二阶风险配置渠道。

影响因素。门控的相对价值不恒定。年度市场波动率最低的三分之一年份，
相对等权的夏普差值平均为 $+0.006$；最高的三分之一年份升至 $+0.189$。
横截面离散度与制度切换次数呈现同向的分组特征，而市场涨跌方向与差值的关系微弱
（秩相关 $-0.049$）。门控的价值来自市场状态的分化程度，而非市场方向。
外生冲击进一步放大了这一分布：把 2015 年股市异常波动、2016 年熔断机制实施、
2018 年中美贸易摩擦、2020 年新冠疫情、2022 年俄乌冲突与疫情封控五个年份
作为冲击年（事件时间由外部事件日历决定），冲击年的夏普差值平均为 $+0.173$，
其余年份为 $+0.073$，差值为 $+0.100$（95\% 区间 $[-0.081,0.276]$）。
冲击年的市场波动率更高（0.291 对 0.202），制度切换次数与其余年份接近（3.20 对 3.29），
门控价值与状态分化的幅度相关，而非切换的频率。

边际效应。单边交易成本每上升 10 个基点，相对等权的夏普差值下降约 0.015
（$p<0.001$），盈亏平衡成本约为 92 个基点。制度状态数呈内点最优：
由 2 增至 3 使夏普差值提高 0.079，由 3 增至 4 转为下降 0.014。
训练窗长度没有单调效应（斜率 $-0.006\,3$/年，$p=0.617$）。

稳健性与边界。门控优势在 2020 年后的子样本中由 $+0.180$ 收窄至 $+0.049$；
四种制度定义下夏普跨度约 0.16，与门控的整体优势同量级；
阈值检验的显著性依赖自助区块长度，3--6 个月时至少一条渠道显著，
12--24 个月时两条渠道均不显著；在 10 只美国 ETF 的窄横截面上互补性不可识别。

相对于以绝对收益为目标、以架构复杂度为创新点的那类多智能体研究，
本文的增量体现在三个方面。理论上，本文把门控价值分解为互补性收益、
估计误差成本与换手成本，给出门控占优的充要条件、两条渠道的分离、
持续性的乘数效应与容量约束下的最优制度数，并把门控收益界定为信息租金，
把「预测组合之谜」解释为该框架的特例。方法上，本文识别了满仓长仓约束下
绝对收益目标的退化问题并给出主动收益的修正，构造了可自校准的互补性阈值检验、
渠道识别检验与 Oracle 上界检验，并用外生冲击年份构成识别设计。
实证上，本文在中国 A 股市场提供严格滚动样本外检验，
用三类独立检验交叉验证门控收益的来源与量级，并给出多智能体策略价值对交易成本、
制度数与训练窗长度的边际效应估计。

余下部分安排：第 \ref{sec:lit} 节综述相关文献；第 \ref{sec:framework} 节给出分析框架；
第 \ref{sec:design} 节说明研究设计；第 \ref{sec:factors} 节报告影响因素；
第 \ref{sec:marginal} 节报告边际效应；第 \ref{sec:robust} 节检验稳健性；
第 \ref{sec:policy} 节讨论政策含义；第 \ref{sec:conclusion} 节总结。
